\section{Introduction}
Vision-language models (VLMs) have emerged as the foundational backbone for multimodal AI systems, enabling autonomous agents to perceive the visual world, reason over multimodal content~\citep{yue2023mmmu}, and interact with both digital~\citep{xie2024osworld,cua2025} and physical environments~\citep{rt2,pi0}. 
The significance of these capabilities has motivated extensive exploration across multiple dimensions, including novel architectural designs~\citep{Alayrac2022FlamingoAV,team2024chameleon,dream2025} and innovative training methodologies with optimized data recipes~\citep{prismaticVLM,nvlm2024}, leading to rapid progress in the field~\citep{liu2023llava,tong2024cambrian,Qwen25VL}.



In this report, we share our efforts to build a compact yet powerful VLM, \mimovl{}. \mimovl{} comprises three components: (1) a native-resolution Vision Transformer (ViT) encoder that preserves fine-grained visual details, (2) a Multi-Layer Perceptron (MLP) projector for efficient cross-modal alignment, and (3) our \mimo{}~\citep{xia2025mimo} language model, specifically optimized for complex reasoning tasks. 

The development of \mimovl{} involves two sequential training processes: 
(1) A \emph{four-stage pre-training phase}, which includes projector warmup, vision-language alignment, general multimodal pre-training, and long-context Supervised Fine-Tuning (SFT). 
Throughout these stages, we curate high-quality datasets by strategically combining open-source collections with synthetic data generation techniques, consuming 2.4 trillion tokens, adjusting the dataset distribution in different stages to facilitate training. This phase yields the \textbf{\mimovlsft{}} model. 
(2) A subsequent \emph{post-training phase}, where we introduce Mixed On-policy Reinforcement Learning (MORL), a novel framework that seamlessly integrates diverse reward signals spanning perception accuracy, visual grounding precision, logical reasoning capabilities, and human preferences. 
We adopt the idea of GRPO~\citep{shao2024deepseekmath} and enhance training stability by exclusively performing on-policy gradient updates during this stage. This phase yields the \textbf{\mimovlrl{}} model.


During this journey, we find:

\textbf{(1) Incorporating high-quality, broad-coverage reasoning data from the pre-training stage is crucial for enhancing model performance. } 
In the current era of ``thinking'' models, vast quantities of multimodal pre-training data are undergoing a significant re-evaluation. Traditional question-answering (QA) data, with its direct, short answers, often restricts models to superficial pattern matching and leads to overfitting. 
In contrast, synthesized reasoning data with long Chain-of-Thought~(CoT) enables learning of complex logical relationships and generalizable reasoning patterns, providing richer supervision signals that markedly improve both performance and training efficiency.
To leverage this advantage, we curate high-quality reasoning data by identifying diverse queries, employing large reasoning models to regenerate responses with long CoT, and applying rejection sampling to ensure quality.
Furthermore, rather than treating this as supplementary fine-tuning data, we incorporate substantial volumes of this synthetic reasoning data directly into the later pre-training stages, where extended training yields continued performance improvements without saturation.


\textbf{(2) Mixed On-policy Reinforcement Learning further enhances model performance, while achieving stable simultaneous improvements remains challenging.}
We apply RL across diverse capabilities, including reasoning, perception, grounding, and human preference alignment, spanning modalities including text, images, and videos.
While this hybrid training approach further unlocks the model's potential, interference across data domains remains challenging.
Disparities in the growth trends of response length and task difficulty levels hinder stable, simultaneous improvements across all capabilities.



\mimovlrl{} delivers exceptional results across the full spectrum of multimodal capabilities.
\textbf{(1) On fundamental visual perception tasks,} it achieves state-of-the-art performance among open-source VLMs of comparable scale, scoring 66.7 on MMMU~\citep{yue2023mmmu} and outperforming Qwen2.5-VL-7B~\citep{Qwen25VL} on 35 out of 40 evaluated tasks. 
\textbf{(2) For complex multimodal reasoning,} \mimovlrl{} exhibits outstanding performance, scoring 59.4 on OlympiadBench~\citep{he2024olympiadbench} and surpassing models up to 72B parameters.
\textbf{(3) In GUI grounding for agentic applications,} our model sets a new standard by achieving a score of 54.7 on OSWorld-G~\citep{osworldg}, even exceeding specialized models like UI-TARS~\citep{qin2025uitars1}.
\textbf{(4) In terms of user experience and preference,} \mimovlrl{} achieves the highest Elo rating among all open-source VLMs on our in-house user preference evaluation, performing competitively with proprietary models such as Claude 3.7 Sonnet.


These results validate our approach: by combining strong perception, sophisticated reasoning, and precise grounding capabilities through our MORL framework, \mimovlsft{} and \mimovlrl{} establish new standards for open-source vision-language models.
To promote transparency and reproducibility, we also contribute a comprehensive evaluation suite covering 50+ tasks with complete prompts and protocols, enabling the community to build upon our work.


